% ============================================================================
% chapitre1.tex — Contexte general du projet
% ============================================================================
\chapter{Contexte general du projet}
\label{chap:contexte}

\introchapitre

Ce premier chapitre situe le projet dans son contexte. Nous presentons d'abord
l'organisme d'accueil, \textbf{HPS Group Morocco}, puis nous etudions le systeme de
detection de fraude actuellement en production afin d'en degager les limites. Nous
exposons ensuite la solution proposee, la methodologie de gestion de projet retenue,
ainsi que le planning et les outils mobilises.

\section{Presentation de l'organisme d'accueil}
\label{sec:organisme}

\textbf{HPS Group} (Hightech Payment Systems) est un editeur et integrateur de
solutions monetiques fonde en 1995 et dont le siege social est situe a Casablanca,
au Maroc. Specialiste des paiements electroniques, HPS Group accompagne banques,
emetteurs, switchs nationaux et commercants a travers sa plateforme phare
\textit{PowerCARD}, qui couvre l'ensemble de la chaine de valeur monetique :
emission, acquisition, switching et gestion des cartes de paiement.

\begin{table}[H]
\centering
\renewcommand{\arraystretch}{1.4}
\begin{tabular}{|l|l|}
\hline
\textbf{Champ} & \textbf{Information} \\
\hline
Raison sociale & HPS Group (Hightech Payment Systems) \\
\hline
Secteur d'activite & Paiements electroniques / Monetique \\
\hline
Siege social & Casablanca, Maroc \\
\hline
Activite principale & Edition de solutions de switching et de routage de paiements \\
\hline
Annee de fondation & 1995 \\
\hline
Effectif & $\sim$500 employes \\
\hline
Produit phare & PowerCARD \\
\hline
\end{tabular}
\caption{Fiche technique de HPS Group}
\label{tab:fiche-hps}
\end{table}

Le present projet a ete realise au sein de l'entite \textbf{HPS Group Morocco}, dans
le cadre d'un stage de fin d'annee oriente vers l'exploration de solutions de traitement
de donnees en temps reel applicables a la detection de fraude bancaire.

\section{Etude de l'existant}
\label{sec:existant}

\subsection{Description de l'existant}

Le systeme de detection de fraude actuellement en production repose sur un
\textbf{pipeline de traitement par lots (batch)} execute chaque nuit. L'ensemble des
transactions de la journee est accumule, puis traite en une seule passe a heure fixe,
produisant des resultats consultables uniquement le lendemain matin. Ce mode de
fonctionnement introduit une \textbf{latence de plusieurs heures} entre l'occurrence
d'une transaction et la detection d'une eventuelle anomalie.

\subsection{Critique de l'existant}

Cette architecture batch presente plusieurs limites majeures, resumees dans le tableau
\ref{tab:limites-existant} :

\begin{table}[H]
\centering
\renewcommand{\arraystretch}{1.3}
\begin{tabularx}{\textwidth}{@{} l X @{}}
\toprule
\textbf{Limite identifiee} & \textbf{Impact concret} \\
\midrule
Latence de plusieurs heures & Fraude detectee bien apres l'autorisation, transaction irreversible \\
\rowcolor{gray!6}
Traitement nocturne unique & Aucune visibilite sur les transactions en cours de journee \\
Absence de reaction immediate & Impossible de bloquer ou d'escalader une transaction suspecte a temps \\
\rowcolor{gray!6}
Accumulation de donnees non analysees & Volume croissant de transactions en attente d'un traitement differe \\
Pas de monitoring temps reel & Incidents operationnels detectes avec retard \\
\bottomrule
\end{tabularx}
\caption{Limites du systeme de traitement existant}
\label{tab:limites-existant}
\end{table}

Dans le domaine des paiements electroniques, ou la fenetre d'action utile face a une
fraude se mesure en secondes, ce delai rend toute intervention corrective largement
inoperante : une transaction frauduleuse autorisee a 14h00 n'est identifiee comme
suspecte que le lendemain, alors que les fonds ont deja ete debites de maniere
definitive.

\subsection{Solution proposee}

Face a ces constats, nous proposons de concevoir un \textbf{pipeline de streaming
temps reel}, deploye sur une infrastructure Kubernetes et reposant exclusivement sur
des technologies \textbf{open source} :

\begin{itemize}
    \item \textbf{Apache Kafka} (distribution Strimzi) pour l'ingestion et le transport
    des evenements de paiement,
    \item \textbf{Apache Flink} (PyFlink) pour le traitement en continu --- decryptage,
    validation, normalisation, optimisation,
    \item \textbf{MinIO} pour la persistance des donnees selon une architecture
    medaillon (Bronze / Silver / Gold),
    \item \textbf{PostgreSQL}, alimente via un connecteur Kafka Connect JDBC Sink,
    pour la persistance finale exploitable par les outils d'analyse,
    \item \textbf{Prometheus} et \textbf{Grafana} pour l'observabilite temps reel de
    l'ensemble du pipeline.
\end{itemize}

Cette solution permet de traiter chaque transaction en quelques secondes, contre
plusieurs heures pour le pipeline batch existant, tout en restant entierement
decouplee de l'infrastructure de production --- garantissant ainsi un PoC sans aucun
impact sur les systemes existants.

\section{Choix du modele de developpement}
\label{sec:scrum}

Nous avons retenu la methodologie \textbf{Scrum} pour la gestion de ce projet. Ce choix
se justifie par la nature exploratoire du PoC : les briques technologiques (Kafka,
Flink, Kafka Connect) necessitaient des iterations courtes pour valider chaque
hypothese technique avant de passer a la suivante. Le travail a ete organise en
sprints courts, rythmes par un backlog priorise et des points d'avancement reguliers
permettant d'ajuster rapidement les choix d'architecture.

\section{Planning previsionnel}
\label{sec:planning}

Le tableau \ref{tab:planning} synthetise le planning previsionnel du sprint principal
de mise en place du pipeline.

\begin{table}[H]
\centering
\renewcommand{\arraystretch}{1.3}
\begin{tabularx}{\textwidth}{@{} l l X @{}}
\toprule
\textbf{Jour} & \textbf{Periode} & \textbf{Objectifs} \\
\midrule
Jour 1 & 20 avril 2026 & Provisionnement Terraform, cluster Kafka (Strimzi), namespaces Kubernetes \\
\rowcolor{gray!6}
Jour 2 & 21 avril 2026 & Developpement des jobs Flink 1 et 2 (decryptage, validation ISO) \\
Jour 3 & 22 avril 2026 & Developpement des jobs Flink 3 et 4 (normalisation, optimisation/scoring) \\
\rowcolor{gray!6}
Jour 4 & 23 avril 2026 & Integration MinIO (medaillon), Kafka Connect, JDBC Sink, PostgreSQL \\
Jour 5 & 24 avril 2026 & Mise en place du monitoring Prometheus/Grafana, tests de bout en bout \\
\bottomrule
\end{tabularx}
\caption{Planning previsionnel du sprint (5 jours, 20--24 avril 2026)}
\label{tab:planning}
\end{table}

\section{Outils et technologies}
\label{sec:outils}

Le tableau \ref{tab:outils-ch1} liste les principaux outils et technologies retenus
pour la realisation de ce PoC, tous open source.

\begin{table}[H]
\centering
\renewcommand{\arraystretch}{1.3}
\begin{tabularx}{\textwidth}{@{} l X @{}}
\toprule
\textbf{Outil / Technologie} & \textbf{Role dans le projet} \\
\midrule
Apache Kafka (Strimzi) & Bus d'evenements distribue, transport des transactions entre etapes \\
\rowcolor{gray!6}
Apache Flink (PyFlink) & Traitement de flux en continu (decryptage, validation, normalisation, scoring) \\
MinIO & Stockage objet S3-compatible, architecture medaillon Bronze/Silver/Gold \\
\rowcolor{gray!6}
PostgreSQL & Base de donnees relationnelle, persistance finale (\texttt{gold\_transactions}) \\
Kafka Connect (Debezium JDBC Sink) & Connecteur assurant le transfert Kafka -> PostgreSQL \\
\rowcolor{gray!6}
Prometheus & Collecte et stockage des metriques du pipeline \\
Grafana & Visualisation et tableaux de bord temps reel \\
\rowcolor{gray!6}
Terraform & Infrastructure as Code, provisionnement reproductible du cluster \\
Kubernetes (Minikube) & Orchestration des conteneurs, environnement de developpement local \\
\rowcolor{gray!6}
Python & Langage principal (producteur, jobs PyFlink, scripts d'integration) \\
Docker & Conteneurisation des composants du pipeline \\
\bottomrule
\end{tabularx}
\caption{Outils et technologies utilises}
\label{tab:outils-ch1}
\end{table}

\conclusionchapitre

Ce chapitre a permis de poser le contexte du projet : un organisme d'accueil reconnu
dans le secteur monetique, un systeme de detection de fraude existant penalise par une
latence batch incompatible avec les exigences du temps reel, et une solution proposee
fondee sur un ecosysteme open source de streaming. Le chapitre suivant approfondit
l'analyse des besoins fonctionnels et non fonctionnels du systeme a concevoir.
